\chapter*{Executive Summary}

This report presents the construction, optimization, and performance evaluation of an equity investment portfolio comprising selected companies listed on the Johannesburg Stock Exchange (JSE). The primary objective was to investigate whether quantitative portfolio optimization techniques could produce a portfolio that delivers superior risk-adjusted performance relative to the Johannesburg Stock Exchange All Share Index (JSE ALSI).

Historical daily price data were collected for seventeen large-cap JSE-listed companies and processed to estimate expected returns, portfolio risk, and asset correlations. Modern Portfolio Theory (MPT) served as the foundation for portfolio construction, with optimization performed under both the Minimum Variance and Maximum Sharpe Ratio criteria. To improve the robustness of portfolio risk estimates, covariance shrinkage techniques were incorporated into the optimization framework.

The optimized portfolio was subsequently evaluated using historical backtesting techniques. Performance was assessed through cumulative returns, wealth growth, portfolio drawdown, and comparison against the JSE ALSI benchmark. The empirical results indicate that the optimized portfolio achieved stronger cumulative performance than the benchmark while maintaining an acceptable level of downside risk over the investment horizon.

Several graphical analyses were produced to support the investment findings, including the efficient frontier, asset correlation heatmap, portfolio allocation charts, drawdown analysis, and benchmark performance comparison. These visualizations provide insight into the diversification characteristics, risk-return trade-off, and overall effectiveness of the optimization strategy.

Overall, the results demonstrate that quantitative portfolio optimization provides a structured and data-driven approach to equity portfolio construction. Although the analysis assumes static portfolio weights and excludes practical considerations such as transaction costs and periodic rebalancing, the framework establishes a strong foundation for more advanced portfolio management applications. Future extensions may incorporate dynamic optimization, alternative risk measures, and machine learning techniques to further enhance portfolio performance.

\clearpage